\section{Empirical strategy and modeling approaches} \label{identification_strategy}

The main goal of our study is to estimate the causal effect of featuring playlists on Spotify's Search Page and compare it to the effect of adding songs by exceptionally popular major label artists to playlists. We first describe the identification challenges associated with these two primary goals and then present two modeling approaches that address these challenges---the two-way fixed effects model and the staggered synthetic difference in difference model. All of our estimates reflect short-term treatment effects as we are relying on the variation in featuring on the Search Page or adding tracks by major label superstars across weeks within a month. Estimates of the long-term effects on playlist followers of these demand factors goes beyond our data coverage.  

\subsection{Identification challenges}
\label{challenges_identification}

The observational nature of the data poses several challenges to identifying treatment effects \citep[e.g.,][]{berman2022value}: (i) every playlist in our data set is either treated in each period (i.e., featured on the Search Page or tracks by popular major label artists being added) or untreated in each period, and we do not observe the counterfactual outcome in the unobserved condition; (ii) the decision of whether to feature a playlist or add popular tracks may be endogenous in the sense that it may be correlated with Spotify's or the focal curators' expectations about changes in demand for the playlist. For example, Spotify may add more of its own playlists to the Search Page to
counteract (amplify) periods of naturally low (high) demand. Relatedly, playlist curators may add tracks by extremely popular major label artists to correct a downward move in followers; (iii) playlists are added to the Search Page or tracks by popular major label artists are added at different times, requiring careful computation and aggregation of treatment effects \citep[e.g., ][]{baker2022much}; %(iv) we do not observe all playlists for all time periods, making it challenging to find a control group;
(iv) Spotify or other curators could have implemented other strategies concurrently when featured on the Search Page or having added ``superstar" tracks, and we might not be able to separate the effect of concurrent changes from featuring or major label track addition; and (v) there may be spillovers between treated and non-treated playlists within the same category, potentially violating the stable unit treatment value assumption \citep[e.g.,][]{goldfarb2022conducting}.  
%Next, we describe two modeling approaches that enable us to estimate the treatment effects of interest. For each model, we highlight---if applicable---the limitations with regard to resolving the challenges for identifying treatment effects, as described before. 


\subsection{Two-way fixed effects model} 

\label{two_way_fe}

The first approach models the main drivers of playlist followers in a reduced-form log-linear demand model, similar to in \cite{shapiro_hitsch_tuchman_2021}. The proposed research design requires fewer assumptions compared to a structural demand model using aggregate data \citep[e.g.,][]{Berry1994} and is suitable for our main goal of comparing the effect of featuring playlists on Spotify's Search Page with the effect of adding tracks by popular major label artists to playlists.\footnote{For example, nested logit models may impose restrictive substitution patterns \citep[e.g.,][]{Berry1995,bresnahan_1996_RAND} and the discrete-choice assumption inherent to structural models of unobserved heterogeneity \citep[][]{Berry1995} is likely violated modeling playlist choice based on followers.} % @Max: Can you please clarify? > "is likely violated modeling playlist choice" (unclear to me)
Moreover, as the proposed model is parsimonious, it allows us to flexibly include characteristics of the main competitor playlists in the same category (leveraging the dynamic categorization derived in Web Appendix \hyperref[details_mining]{C}). We specify the two-way fixed effects model as follows:

\begin{align} 
\begin{split} 
\text{log}(f_{pt}+k) =
& \gamma_{p,m(t)} + \gamma_t + \eta^{\text{feat}} \text{Treated}_{pt}^{\text{feat}} + \eta^{\text{star}} \text{Treated}_{pt}^{\text{star}} + \\ & + \beta \text{log}(ml_{pt}+k) + \omega \text{log}(r_{C(p),t}+k) + \phi^{\prime} x_{pt} + \psi^{\prime} x_{-pt} + \varepsilon_{pt}, 
\end{split}
\label{two_way_model}
\end{align} %\text{log}(1+n_{pt})

where $f_{pt}$ measures the followers of playlist $p$ on day $t$.\footnote{We set the constant to a small number, $k=1$. In a robustness check, we also consider the Inverse Hyperbolic Sine (IHS) transformation to deal with zero-valued observations \citep{ihs}.} The coefficients $\eta^{feat}$ and $\eta^{star}$ measure the effect of including playlists on Spotify's Search Page as well as the effect of adding at least one track by a highly popular major label artist (``superstar") to playlists on day $t$, where $\text{Treated}_{pt}^{\text{feat}}, \text{Treated}_{pt}^{\text{star}} \in \{0,1\}$. %\footnote{Recall that we define a superstar as any artist appearing in the Spotify Top 10 between December 2016 and July 2019, or the US Billboard Top 10 between 1999 and 2019.} %Hannes: investigate how long Spotifycharts has been around - maybe 2016 is start of available data so we should mention that.
Moreover, we include a rich set of (time-varying) fixed effects and $\gamma_{p,m(t)}$ denote playlist-month fixed effects, i.e., $m(t) \in \{M_1, M_2, \dots \}$. In addition, $\gamma_t$ are date-fixed effects, accounting for common demand factors across playlists within a date. 

The second line in Equation \ref{two_way_model} accounts for various control terms. $\beta$ controls for the effect of including major label tracks on the playlist, beyond those from exceptionally popular superstars, where $ml_{pt}$ measures the share of major label tracks on playlist $p$. $\omega$ measures the effect of personalized recommendations in category $C(p)$ at date $t$ on the followers of playlist $p$.\footnote{$r_{C(p),t}$ reflects the ``recommendation intensity", defined as the number of songs added to personalized playlists in category $C(p)$ at date $t$, divided by the number of personalized playlists in category $C(p)$.
Among the user-curated playlists in the raw data, we observe a sample of $1,634$ of Spotify's personalized playlists, such as ``Discover Weekly" (60\% of all personalized playlists), users' compilations of the most listened-to songs from previous years ($31\%$) and ``Release Radar" ($6\%$). Personalized playlists are typically private to individual users, but can become public when shared on social media.}
%We observe users' ``Discover Weekly" playlists (60\% of all personalized playlists), users' compilation of the most listened-to songs from previous years (31\%), ``Release Radar" (6\%) and other remaining personalized playlists (e.g., “Your Summer Rewind”, 1\%; “Daily Mix”, 1\%).
$\phi^{\prime} x_{pt}$ controls for the average track age, playlist duration (in hours) as well as acoustic
attributes of playlist $p$, while $\psi^{\prime} x_{-pt}$ controls for featuring on the Search Page and availability of superstar tracks of the fifteen largest competing playlists in the same category. %\footnote{We collapse the smallest $5\%$ of categories in terms of followers as a composite rest category, ensuring that we always observe at least 15 competitors in a category.}  
%obtained from the categorization using association rule mining (Web Appendix \hyperref[details_mining]{C}). 
Including the rich control variables in Equation \ref{two_way_model} is important to address identification concern (iv) and ensure that the estimated treatment effects are not confounded with simultaneously occurring changes in other playlist covariates.

Figure \ref{fig:resid_var} in Web Appendix \hyperref[appendix:appendix_model]{D} shows the residual identification variation in featuring (left panel) and additions of tracks by major label superstars (right panel) after accounting for playlist-month fixed effects and remaining control variables as in Equation \ref{two_way_model}. There is sufficient variation in both variables left that is caused by Spotify updating its Search Page (see Figure \ref{fig:section_updates}) and playlist curators adding tracks by major label superstar artists within weeks of a month.\footnote{The median update frequency for an average playlist in our sample equals 14 days.}

The two-way fixed effects model assumes that conditional on the rich fixed effects and control variables, the residual variation in Figure \ref{fig:resid_var} is quasi-random---i.e., not correlated with unobserved demand factors. To justify this assumption, one might argue that user demand for a playlist does not drastically change within a month. Hence, variations in the composition of the playlist's songs and the Search Page are quasi-random within weeks of a month. However, we cannot observe the stability of users' preferences for playlists, and Spotify may update its playlists/Search Page in response to demand shocks within a month. %, but about 54\% of playlists update their content at a higher frequency than the weekly level.} 
In other words, the main limitation of this model pertains to the identification challenge (ii) as there might be endogenous selection of remaining playlists into the treatment and control groups, even conditional on including playlist-month fixed effects. To address this concern, we would ideally want to create a control group of playlists that closely resemble treated playlists in the days before the event but do not get treated in our data.\footnote{An alternative strategy would be to use instrumental variables to resolve endogeneity concerns. However, it is very difficult to identify valid instruments for featuring playlists and including popular superstar tracks in playlists that fulfill the exclusion restriction assumption.} Another potential issue with the two-way fixed effects model is that it does not adequately account for the different timing of treatment effects (challenge iii), potentially biasing estimation results if treatment effects can change over time (``dynamic treatment effects") due to a ``bad comparisons" problem \citep{baker2022much}. To mitigate these limitations, we propose an alternative model. 

%%%%%%%%%%%%%%%%%%%%%%%%%%%%%%%%%%%%%%%%%%%
%%%%%%%%%%%%%%%%%%%%%%%%%%%%%%%%%%%%%%%%%%%
%%%%%%%%%%%%%%%%%%%%%%%%%%%%%%%%%%%%%%%%%%%
\subsection{Staggered synthetic difference in difference (SDID)} 

\label{SDID}
%%%%%%%%%%%%%%%%%%%%%%%%%%%%%%%%%%%%%%%%%%%
%%%%%%%%%%%%%%%%%%%%%%%%%%%%%%%%%%%%%%%%%%%
%%%%%%%%%%%%%%%%%%%%%%%%%%%%%%%%%%%%%%%%%%%

To investigate whether---and if so, by how much---the two-way fixed effects specification (Equation \ref{two_way_model}) biases treatment effect estimates, we use the synthetic difference in difference model \citep{SDiD_AER_Arkhangelsky_21}. By weighting the observations of never-treated playlists to generate ``clones" of the treated observations, this method improves upon our identification strategy about the endogenous selection into treatment and control groups (challenge ii). Moreover, we account for identification challenge (iii)---the differential treatment timing---by repeatedly applying the SDID estimator on the subset of data implied by unique adoption dates, as recommended in \cite{SDiD_AER_Arkhangelsky_21}.

The SDID's crucial ``exclusion restriction" assumption is that the synthetic control playlists precisely match the follower evolution of the treated playlists in the times before the event.\footnote{The SDID method compensates for the lack of parallel trends between treated and control playlists by reweighting units to match their pre-exposure trends. This procedure weakens the reliance on parallel trend assumptions as in conventional difference in difference methods \citep{SDiD_AER_Arkhangelsky_21}.} Treating a playlist can then be regarded as quasi-random compared to its synthetic control unit. To justify this assumption, we present the Mean Squared Percentage Error (MSPE) between the treated playlist's followers and the synthetic control playlists over the 15 days prior to the event, detailed in Table \ref{tab:model_results}. The small MSPE highlights a strong fit, facilitated by leveraging thousands of never-treated donor units to generate synthetic ``clones" for treated playlists. There are two potential challenges (iv and v) to this exclusion restriction assumption that we address as follows. First, Spotify or other curators may have concurrently implemented other changes after treatment (identification challenge iv). We checked the stability in playlists' characteristics after treatment in Table \ref{tab:covbalance} in Web Appendix \hyperref[appendix:stability_test]{E}. We found that curators changed almost none of the covariates after their playlists got featured on the Search Page or tracks by exceptionally popular major label artists were added. %Importantly, Table \ref{tab:covbalance} confirms that playlists that get included on the Search Page do not add major label superstars in response. Similarly, playlists that added tracks by major label superstars do not get featured with a higher likelihood.      
Second, we use only donor playlists from categories other than the treated units to ensure that the stable unit treatment value assumption is not violated \citep[][]{goldfarb2022conducting}.   

%We now explain how we operationalize estimating the staggered difference in difference model. We first outline how we select events in the two treatment cases---featuring on Spotify's Search Page and adding tracks by highly popular major label artists to playlists---and discuss the econometric specification afterward. 

%\subsubsection{Event selection}
%\label{event_selections}
\subsubsection{Playlist featuring} 
%\textbf{\textit{Playlist featuring.}} 
We use data on all playlists we observed being featured for the first time. To have sufficient pre-treatment data, we require playlists to have been observed not featured for at least two weeks. We retain observations 15 days before and 15 days after treatment and require that playlists remain treated in the after-period for at least seven days. Choosing such a narrow analysis window mitigates the impact of confounding variables \citep[e.g.,][]{orhun_16}. These steps led to a sample of 113 treated playlists on 23 unique dates. 

Featured observations include more popular playlists such as ``Reggaeton Classics" ($1.9$ million followers) as well as less popular playlists such as  ``Tech House Chargers" ($4.5$k followers). Most of the featured playlists were curated by Spotify, indicating that Spotify tends to prefer selecting its own playlists for inclusion on the Search Page.\footnote{Initially, about 15\% of playlists that Spotify featured were from other curators. At the end of November 2019, this share dropped to virtually zero, where it remained for the rest of our sample.} Figure \ref{treatmentvar_promo} in Web Appendix \hyperref[appendix:appendix_model]{D} illustrates the variation in treatment dates across treated playlists. In line with Figure \ref{fig:section_updates}, we see that playlists were included on many different dates during our observation period.

\subsubsection{Addition of tracks from popular major label artists}
We define the treatment as adding at least one track to a playlist by a major label ``superstar" artist included in the Spotify Top 10 between December 2016 and July 2019, or the US Billboard Top 10 between 1999 and
2019. Artists include established superstars such as Alicia Keys, Jennifer Lopez, and Ed Sheeran, as well as more recent superstars such as Shawn Mendes, Billie Eilish, and Lizzo. Again, we keep treated playlists that are observed at least 15 days before and 15 days after treatment, with the requirement that newly added tracks be listed for at least seven days after the event. It is possible that playlists can be ``treated" multiple times; in these few cases, we select only the first treated event of a playlist (under the restriction that the second treatment did not happen in the fifteen days after the first treatment). These steps lead to a sample of 2,248 treated playlists on 107 days. Figure \ref{treatmentvar_ml} in Web Appendix \hyperref[appendix:appendix_model]{D} shows the treatment variation plot for the selected events.
% clarify "it is possible that playlists can get" % 
% max says: we need the table


\subsubsection{Econometric specification}
Conditional on every treatment date $a \in A$, we have a balanced panel with $P_{co,a}$ never-treated control playlists and $P_{tr,a}$ treated playlists that get exposed after $T_{pre} = 15$ and are observed for $T_{post} = 15$ consecutive time periods, with $P_a=P_{co,a}+P_{tr,a}$ and $T=T_{pre}+T_{post}$. We follow \cite{SDiD_AER_Arkhangelsky_21} and estimate the average causal effect of exposure for batch $a$ in a two-way fixed effects model:

\begin{align} 
\begin{split} 
\left(\hat{\eta}_a^{\text{con}}, \hat{\mu}_a, \hat{\gamma}_{p,a}, \hat{\gamma}_{t,a}\right) = \underset{\eta^{\text{con}},\mu, \gamma_{p}, \gamma_{t}}{\text{arg min}} \left\{\sum_{p=1}^{P_a} \sum_{t=1}^T \left( \tilde{f}_{pt} - \mu - \gamma_{p} - \gamma_{t} - \eta^{\text{con}} \text{Treated}_{pt}^{\text{con}} \right)^2 \hat{\omega}_{p,a} \hat{\lambda}_{t,a} \right\}, \end{split} 
\label{SDID_speci}
\end{align} 

where $\tilde{f}_{pt} = \text{log}(f_{pt}+k)$, con $\in \{\text{feat},\text{star}\}$, $\hat{\omega}_{p,a}$ are weights aligning pre-exposure trends in log-followers of never-treated units with those of treated units, $\hat{\lambda}_{t,a}$ are time weights balancing pre-exposure time periods with postexposure ones.\footnote{\cite{SDiD_AER_Arkhangelsky_21} show that the inclusion of unit and time weights makes the estimator more robust and improves precision.} Because of our staggered adoption setting, we apply Algorithm 1 in \cite{SDiD_AER_Arkhangelsky_21} on every of the $|A|$ balanced panel data sets. We compute the aggregate treatment effect as the average effect across all event dates weighted by the event date-specific number of treated playlists \citep{SDiD_AER_Arkhangelsky_21}.\footnote{For each adoption date $a \in A$, the variance of the average treatment effect estimator is computed using the placebo method \citep[Algorithm 4 in][]{SDiD_AER_Arkhangelsky_21} using 200 replications. Global standard errors are given as the square root of the variance of the aggregate treatment effect estimator \citep[e.g.,][]{Imbens_JoEconmetrics_22}.}%, i.e., $SE^{\text{con}} = \sqrt{\text{Var}\left(\sum_a w_a \hat{\eta}_a^{\text{con}}\right)} = \sqrt{\sum_a(w_a^2*\text{Var}(\hat{\eta}^{\text{con}}_a))}$, where $w_a= P_{tr,a}/(\sum_a P_{tr,a})$ are the adoption date-specific weights.}  

%Similar to Berman and Israeli (2022) and Li et al. (2022), we compute the average treatment effect across all cohorts (), as well as its standard error, by weighting the cohort-specific effects by the respective number of treated households (Arkhangelsky et al. 2021, p. 4114). 

%A potential concern with the synthetic difference in difference model in Equation \ref{SDID_speci} is that Spotify or other curators may have concurrently implemented other changes after treatment (identification challenge iv). We checked the stability in playlists' characteristics after treatment in Table \ref{tab:covbalance} in Web Appendix \hyperref[appendix:stability_test]{F}. We found that curators changed almost none of the covariates after their playlists got featured on the Search Page or tracks by exceptionally popular major label artists were added. Importantly, Table \ref{tab:covbalance} confirms that playlists that get included on the Search Page do not add major label superstars in response. Similarly, playlists that added tracks by major label superstars do not get featured with a higher likelihood.   


